\apendice{Documentación técnica de programación}

\section{Introducción}

Este anexo describe la estructura técnica del repositorio y los pasos para ejecutar, probar y desplegar la solución.

\section{Estructura de directorios}

\begin{itemize}
    \item \texttt{src/app.py}: API principal Flask.
    \item \texttt{src/classifier.py}: clasificador automático de incidencias.
    \item \texttt{src/storage.py}: persistencia local o en Azure Blob Storage.
    \item \texttt{src/config.py}: configuración por variables de entorno.
    \item \texttt{src/requirements.txt}: dependencias Python.
    \item \texttt{tests/test\_api.py}: pruebas unitarias.
    \item \texttt{infra/terraform/}: infraestructura como código.
    \item \texttt{logicapp/}: definición de Logic App.
    \item \texttt{.github/workflows/}: CI/CD con GitHub Actions.
    \item \texttt{memoria/tex/}: memoria y anexos en LaTeX.
\end{itemize}

\section{Ejecución local}

\begin{verbatim}
cd src
pip install -r requirements.txt
set STORAGE_MODE=local
python app.py
\end{verbatim}

La API quedará disponible en \texttt{http://localhost:5000}.

\section{Variables de entorno}

\begin{itemize}
    \item \texttt{STORAGE\_MODE}: \texttt{local} o \texttt{azure}.
    \item \texttt{AZURE\_STORAGE\_CONNECTION\_STRING}: cadena de conexión del Storage Account.
    \item \texttt{KEY\_VAULT\_URL}: URI del Key Vault.
    \item \texttt{API\_KEY}: token Bearer opcional para desarrollo local.
\end{itemize}

\section{Despliegue en Azure}

\begin{enumerate}
    \item Configurar el secreto \texttt{AZURE\_CREDENTIALS} en GitHub.
    \item Ejecutar \texttt{terraform apply} en \texttt{infra/terraform}.
    \item Desplegar la carpeta \texttt{src} en App Service mediante GitHub Actions.
    \item Importar la Logic App desde \texttt{logicapp/logicapp-workflow.json}.
\end{enumerate}

\section{Pruebas del sistema}

\begin{verbatim}
python -m unittest discover -s tests -p "test_*.py"
\end{verbatim}

Las pruebas verifican endpoints, creación de incidencias, métricas y clasificación automática.